\section{Preliminary}

\(D\)차원 벡터 \(\mathbf{x}_t \in \mathbb{R}^D\)가 이산 시간 \(t\)에서의 K-line 관측치를 나타낸다고 하자. 이는 \(D\)개의 핵심 금융 지표로 구성된다. 본 연구에서는 OHLCVA 속성(Open, High, Low, Close 가격, trading Volume, Amount)을 나타내기 위해 차원을 \(D=6\)으로 고정한다. 이러한 입력 선택의 근거는 Appendix~\ref{sec:discussion} (Q1)에 자세히 설명되어 있다. 과거 시퀀스
\(
\mathbf{x}_{1:T} = (\mathbf{x}_1, \mathbf{x}_2, \ldots, \mathbf{x}_T),
\)
가 주어졌을 때, 우리의 목표는 다음 \(H\)개의 관측치
\(
\hat{\mathbf{x}}_{T+1:T+H} = (\hat{\mathbf{x}}_{T+1}, \hat{\mathbf{x}}_{T+2}, \ldots, \hat{\mathbf{x}}_{T+H}).
\)
를 예측하는 것이다.

Kronos는 원시 연속 입력에 직접 작동하는 대신, 먼저 각 다변량 관측치 \(\mathbf{x}_t\)를 이산 토큰 \(b_t\)로, 학습 가능한 코드북 \(\mathcal{C}\)를 통해 양자화한다.
% 이 접근법은 주어진 시점의 전체 OHLCVA 지표 집합을 양자화를 위한 단일 단위로 취급한다.
따라서 원래 시퀀스
\(
\mathbf{x}_{1:T} = (\mathbf{x}_1, \dots, \mathbf{x}_T)
\)
는
\(
\mathbf{b}_{1:T} = (b_1, \dots, b_T).
\)
로 매핑된다. 그러면 예측 과제는 자기회귀 토큰-시퀀스 모델링 문제로 환원된다:
\begin{equation}
\label{eq:ts_ar}
p(\mathbf{b}_{T+1:T+H} \mid \mathbf{b}_{1:T})
= \prod_{h=1}^{H} p\bigl(b_{T+h} \mid \mathbf{b}_{1:T+h-1}\bigr).
\end{equation}

이러한 이산적 정식화는 본질적으로 확장 가능하며, 합성 데이터 생성 및 변동성 예측과 같이 생성적으로 틀 지을 수 있는 다른 과제로 자연스럽게 확장된다.
